%%%==============EX_1============%%%
\begin{ex}%[2H5C1-5]
	Trong không gian với hệ tọa độ $Oxyz$, cho mặt phẳng $(P)\colon x-2y+2z-3=0$ và mặt cầu $(S)\colon x^2+y^2+z^2+2x-4y-2z+5=0$. Giả sử $M\in(P)$ và $N\in(S)$ sao cho $\overrightarrow{MN}$ cùng phương với vectơ $\vec{u}=(1; 0; 1)$ và khoảng cách giữa $M$ và $N$ lớn nhất. Tính $MN$.
	\choice
	{$MN=3$}
	{$MN=1+2 \sqrt{2}$}
	{\True $MN=3 \sqrt{2}$}
	{$MN=14$}
	\loigiai{
		\begin{center}
			\begin{tikzpicture}[>=stealth,line join=round,line cap=round,font=\footnotesize,scale=1]
				\def \r{2};
				\path
				(0,0) coordinate (I)
				(90:\r) coordinate (N')
				(20:\r) coordinate (N)
				($(N')!2.5!(I)$) coordinate (H')
				($(H')!1.5!90:(I)$) coordinate (M)
				($(M)!(N)!(H')$) coordinate (H)
				;
				\draw
				(M)--(H)--(N)--cycle (N')--(H')
				(I) circle (\r)
				;
				\foreach \t/\g in {M/170,N/40,H/-70,H'/-70,I/60,N'/90}{\draw[fill=black] (\t) circle (1pt) node[shift={(\g:7pt)},font=\scriptsize]{$ \t $};}
			\end{tikzpicture}
		\end{center}
		Mặt phẳng $(P)$ có vtpt $\vec{n}=(1;-2; 2)$. \\
		Mặt cầu $(S)$ có tâm $I(-1; 2; 1)$ và bán kính $r=1$. \\
		Nhận thấy rằng góc giữa $\vec{u}$ và $\vec{n}$ bằng $45^{\circ}$. \\
		Vì $d(I,(P))=2> 1=r$ nên $(P)$ không cắt $(S)$.\\
		Gọi $H$ là hình chiếu của $N$ lên $(P)$ thì $\widehat{NMH}=45^{\circ}$ và $MN=\dfrac{NH}{\sin 45^{\circ}}=NH \sqrt{2}$ nên $MN$ lớn nhất khi và chỉ khi $NH$ lớn nhất. Điều này xảy ra khi $N\equiv N^{\prime}$ và $H\equiv H^{\prime}$ với $N^{\prime}$ là giao điểm của đường thẳng $d$ qua $I$, vuông góc $(P)$ và $H^{\prime}$ là hình chiếu của $I$ lên $(P)$.\\
		Lúc đó $NH_{\max}=N^{\prime} H^{\prime}=r+d(I,(P))=3$ và $MN_{\max}=\dfrac{NH_{\max}}{\sin 45^{\circ}}=3 \sqrt{2}$.
	}
\end{ex}
%%%==============EX_2============%%%
\begin{ex}%[2H5C1-3]
	Cho $A(0; 8; 2)$ và mặt cầu $(S)\colon (x-5)^2+(y+3)^2+(z-7)^2=72$ và điểm $A(9;-7; 23)$. Viết phương trình mặt phẳng $(P)$ đi qua $A$ và tiếp xúc với mặt cầu $(S)$ sao cho khoảng cách từ $B$ đến mặt phẳng $(P)$ là lớn nhất. Giả sử $\vec{n}=(1; m; n)$ là một vectơ pháp tuyến của $(P)$. Lúc đó
	\choice
	{$m\cdot n=4$}
	{$m\cdot n=2$}
	{\True $m\cdot n=-4$}
	{$m\cdot n=-2$}
	\loigiai{
		$(P)$ đi qua điểm $A(0;8;2)$ và có vectơ pháp tuyến $\vec{n}=(1;m;n) \Rightarrow(P)\colon x+my+nz-8m-2n=0$.\\
		$(P)$ tiếp xúc với mặt cầu $(S) \Rightarrow \dfrac{|5-11m+5n|}{\sqrt{1+m^2+n^2}}=6\sqrt{2}$.\\
		$\mathrm{d}=\mathrm{d}(B;(P))=\dfrac{|9-15m+21n|}{\sqrt{1+m^2+n^2}}=\dfrac{|5-11m+5n+4-4m+16n|}{\sqrt{1+m^2+n^2}}$\\
		$\le \dfrac{|5-11m+5n|}{\sqrt{1+m^2+n^2}}+4\dfrac{|1-m+4n|}{\sqrt{1+m^2+n^2}}$.\\
		$\leq 6\sqrt{2}+4\dfrac{\sqrt{1^2+(-1)^2+4^2} \cdot \sqrt{1+m^2+n^2}}{\sqrt{1+m^2+n^2}}$ (Buinhiacôpxki).\\
		$=18\sqrt{2}$.\\
		$\Rightarrow \mathrm{d}_{\max}=18\sqrt{2} \Leftrightarrow \dfrac{1}{1}=\dfrac{-1}{m}=\dfrac{4}{n} \Rightarrow\heva{&m=-1\\&n=4} \Rightarrow m \cdot n=-4$.
	}
\end{ex}
%%%%==============EX_3============%%%
\begin{ex}%[2H5C3-3]
	Trong không gian $Oxyz$, cho mặt cầu $(S)\colon(x-3)^2+(y-1)^2+z^2=4$ và đường thẳng
	$d\colon\heva{&x=1+2t \\ &y=-1+t\\& z=-t},(t \in \mathbb{R})$. Mặt phẳng chứa $d$ và cắt $(S)$ theo một đường tròn có bán kính nhỏ nhất có  phương trình là
	\choice
	{\True $y+z+1=0$}
	{$x+3y+5z+2=0$}
	{$x-2y-3=0$}
	{$3x-2y-4z-8=0$}
	\loigiai{
		Mặt cầu $(S)$ có tâm $I(3;1;0)$ và bán kính $R=2$.\\
		Gọi $H$ là hình chiếu của $I$ trên $d$.\\
		$H\in d \Leftrightarrow H(1+2t;-1+t;-t); \overrightarrow{IH}=(-2+2t;-2+t;-t)$.\\
		Véctơ chỉ phương của $d$ là $\overrightarrow{u}_d=(2; 1;-1)$.\\
		$\overrightarrow{IH}\cdot \overrightarrow{u}_d=0\Leftrightarrow 2(-2+2t)+1(-2+t)+t=0\Leftrightarrow t=1$.\\
		 Do đó $H(3; 0;-1) \Rightarrow IH=\sqrt{2}$.\\
		Gọi $(P)$ là mặt phẳng chứa đường thẳng $d$ và cắt mặt cầu $(S)$ theo đường tròn có bán kính $r$. Ta có
		$r=\sqrt{R^2-[\mathrm{d}(I,(P))]^2}=\sqrt{4-[\mathrm{d}(I,(P))]^2}$.\\
		Mà $\mathrm{d}(I,(P)) \le IH=\sqrt{2}$ nên $r=\sqrt{4-[\mathrm{d}(I,(P))]^2} \geq \sqrt{4-IH^2}=\sqrt{4-(\sqrt{2})^2}=\sqrt{2}$.\\
		Suy ra $\min r=\sqrt{2}$, đạt được khi $IH\perp(P)$.\\
		Khi đó mặt phẳng $(P)$ đi qua $H(3; 0;-1)$ nhận $\overrightarrow{IH}=(0;-1;-1)$ làm một véctơ pháp tuyến.\\
		Phương trình mặt phẳng $(P)$ là $0(x-3)-1(y-0)-1(z+1)=0 \Leftrightarrow y+z+1=0$.
	}
\end{ex}
%%%%==============EX_4============%%%
\begin{ex}%[2H5C3-3]
	Trong không gian với hệ tọa độ $Oxyz$, cho hai điểm $A(3;-2;6)$, $B(0;1;0)$ và mặt cầu $(S)\colon (x-1)^2+(y-2)^2+(z-3)^2=25$. Mặt phẳng $(P)\colon a x+b y+c z-2=0$ đi qua $A$, $B$ và cắt $(S)$ theo giao tuyến là đường tròn có bán kính nhỏ nhất. Tính $T=a+b+c$.
	\choice
	{\True $T=3$}
	{$T=4$}
	{$T=5$}
	{$T=2$}
	\loigiai{
		\begin{center}
			\begin{tikzpicture}[>=stealth,line join=round,line cap=round,font=\footnotesize,scale=1]
				\def \r{3};
				\path
				(0,0) coordinate (I)
				(-30:\r) coordinate (C)
				($(I)-(0,1.5)$) coordinate (H)
				($(H)+(1,0)$) coordinate (K)
				($(H)+(1.3,0.25)$) coordinate (B)
				($(B)!5!(K)$) coordinate (A)
				($(A)!0.5!(K)$) coordinate (M)
				;
				\draw (I) circle (\r)
				(C) arc (0:-180: 2.6 cm and 0.4 cm)
				(A)--(M) 
				;
				\draw [dashed]
				(C) arc (0:180: 2.6 cm and 0.4 cm)
				(I)--(H)--(K)--(B)--(M) (I)--(-150:\r)--(H)
				;
				\foreach \t/\g in {A/0,I/90,H/-90,K/150,B/-30}{\draw[fill=black] (\t) circle (1pt) node[shift={(\g:7pt)},font=\scriptsize]{$ \t $};}
			\end{tikzpicture}
		\end{center}
		Mặt cầu $(S)$ có tâm $I(1;2;3)$ và bán kính $R=5$.\\
		Ta có $\heva{&A\in(P)\\&B\in(P)}\Leftrightarrow\heva{&3a-2b+6c-2=0\\ &b-2=0} \Leftrightarrow\heva{&a=2-2c \\ &b=2}$.\\
		Bán kính của đường tròn giao tuyến là $r=\sqrt{R^2-[\mathrm{d}(I,(P))]^2}=\sqrt{25-[\mathrm{d}(I,(P))]^2}$.\\
		Bán kính của đường tròn giao tuyến nhỏ nhất khi và chỉ khi $\mathrm{d}(I;(P))$ lớn nhất.\\
		Ta có $\mathrm{d}(I,(P))=\dfrac{|a+2b+3c-2|}{\sqrt{a^2+b^2+c^2}}=\dfrac{|2-2c+4+3c-2|}{\sqrt{(2-2c)^2+2^2+c^2}}=\sqrt{\dfrac{(c+4)^2}{5c^2-8c+8}}$.\\
		Xét $f(c)=\sqrt{\dfrac{(c+4)^2}{5c^2-8c+8}} \Rightarrow f'(c)=\dfrac{-48c^2-144c+192}{\left(5c^2-8c+8\right)^2 \sqrt{\dfrac{(c+4)^2}{5c^2-8c+8}}}$.\\
		$f'(c)=0\Leftrightarrow\hoac{&c=1\\ &c=-4.}$\\
		Bảng biến thiên
\begin{center}
	\begin{tikzpicture}
		\tkzTabInit[nocadre=false,lgt=2,espcl=2,deltacl=0.5]
		{$c$/1 , $f'(c)$/1 , $f(c)$/2}
		{$-\infty$,$-4$,$1$,$+\infty$}
		\tkzTabLine{,-,0,+,0,-,}
		\tkzTabVar{+/$\dfrac{1}{\sqrt{5}}$,-/$0$,+/$\sqrt{5}$,-/$\dfrac{1}{\sqrt{5}}$}
	\end{tikzpicture}
\end{center}
		Vậy $\mathrm{d}(I,(P))$ lớn nhất bằng $\sqrt{5}$ khi và chỉ khi $c=1\Rightarrow a=0,\, b=2\Rightarrow a+b+c=3$.
	}
\end{ex}
%%%%==============EX_5============%%%
\begin{ex}%[2H5C3-3]
	Trong không gian $Oxyz$, cho hai điểm $A(1;2;4)$, $B(0;0;1)$ và mặt cầu $(S)\colon (x+1)^2+(y-1)^2+z^2=4$. Mặt phẳng $(P)\colon ax+by+cz-4=0$ đi qua $A$, $B$ và cắt $(S)$ theo giao tuyến là một đường tròn có bán kính nhỏ nhất. Tính $T=a+b+c$?
	\choice
	{$T=\dfrac{1}{5}$}
	{$T=\dfrac{3}{4}$}
	{\True $T=1$}
	{$T=-2$}
	\loigiai{
		Ta có $(S)$ có tâm $I(-1;1;0)$ và bán kính $R=2$.\\
		Do $A, B\in(P) \Leftrightarrow\heva{&a+2b+4c-4=0\\ &c-4=0} \Leftrightarrow\heva{&a=-2b-12\\ &c=4.}$\\
		$\Rightarrow(P)\colon -2(b+6) x+b y+4z-4=0$.\\
		Gọi $r$ là bán kính của đường tròn là giao tuyến của $(P)$ và $(S) \Rightarrow r=\sqrt{R^2-\mathrm{d}^2(I,(P))}$, để $r$ đạt giá trị nhỏ nhất $\Leftrightarrow \mathrm{d}(I,(P))$ đạt giá trị lớn nhất. \\
		Mà $\mathrm{d}(I,(P))=\dfrac{|3b+8|}{\sqrt{5b^2+48b+160}}$.\\
		Xét hàm số $f(x)=\dfrac{3x+8}{\sqrt{5x^2+48x+160}}$.\\ $f'(x)=\dfrac{32x+288}{\left(\sqrt{5x^2+48x+160}\right)^3}$.\\ $f'(x)=0\Leftrightarrow x=-9$.\\
		Bảng xét biến thiên
\begin{center}
	\begin{tikzpicture}
		\tkzTabInit[nocadre=false,lgt=2,espcl=4,deltacl=0.5]
		{$x$/1 , $f'(x)$/1 , $f(x)$/2}
		{$-\infty$,$-9$,$+\infty$}
		\tkzTabLine{,-,0,+,}
		\tkzTabVar{+/$-1{,}34$,-/$f(-9)\approx -1{,}66$,+/$1{,}34$}
	\end{tikzpicture}
\end{center}
	Từ đó ta suy ra bảng biến thiên của hàm số $y=|f(x)|$ 
	\begin{center}
		\begin{tikzpicture}
			\tkzTabInit[nocadre=false,lgt=2,espcl=4,deltacl=0.5]
			{$x$/1, $f(x)$/2}
			{$-\infty$,$-9$,$-\dfrac{8}{3}$,$+\infty$}
			\tkzTabVar{-/$1{,}34$,+/$ 1{,}66$,-/$0$,+/$1{,}34$}
		\end{tikzpicture}
	\end{center}
		Dựa vào bảng biến thiên, ta có $x=-9 \Rightarrow b=-9 \Rightarrow a=6 \Rightarrow T=1$.
	}
\end{ex}
%%%%==============EX_6============%%%
\begin{ex}%[2H5C3-3]
	Trong không gian với hệ toạ độ $Oxyz$, cho mặt cầu $(S)\colon (x-1)^2+(y-2)^2+(z-3)^2=9$, điểm $A(0;0;2)$. Mặt phẳng $(P)$ qua $A$ và cắt mặt cầu $(S)$ theo thiết diện là hình tròn $(C)$ có diện tích nhỏ nhất, phương trình $(P)$ là
	\choice
	{$(P)\colon x-2y+3z-6=0$}
	{$(P)\colon x+2y+3z-6=0$}
	{$(P)\colon 3x+2y+2z-4=0$}
	{\True $(P)\colon x+2y+z-2=0$}
		\loigiai{
	\begin{center}
		\begin{tikzpicture}[>=stealth,line join=round,line cap=round,font=\footnotesize,scale=1]
			\def \r{3};
			\path
			(0,0) coordinate (I)
			(0:\r) coordinate (B)
			(-45:\r) coordinate (C)
			($(I)-(0,2.12)$) coordinate (H)
			(-135:\r) coordinate (A)
			;
			\draw (I) circle (\r)
			(B) arc (0:-180: \r cm and 0.8 cm)
			(C) arc (0:-180: 2.12 cm and 0.25 cm)
			;
			\draw [dashed]
			(B) arc (0:180: \r cm and 0.8 cm)
			(C) arc (0:180: 2.12 cm and 0.25 cm)
			(I)--(H)--(A)--cycle
			;
			\foreach \t/\g in {I/90,H/0,A/180}{\draw[fill=black] (\t) circle (1pt) node[shift={(\g:7pt)},font=\scriptsize]{$ \t $};}
		\end{tikzpicture}
	\end{center}
		Mặt cầu $(S)$ có tâm $I(1; 2; 3)$, bán kính $R=3$.\\
		Ta có $IA=\sqrt{6} < R\Rightarrow A$ nằm trong mặt cầu $(S)$.\\
		Do đó mặt phẳng $(P)$ qua $A$ luôn cắt mặt cầu $(S)$ theo thiết diện là hình tròn $(C)$ có bán kính $r=\sqrt{R^2-IH^2}$ (với $H$ là hình chiếu của $I(1;2; 3)$ trên $(P)$).\\
		Ta luôn có $IA\geq IH\Rightarrow \sqrt{R^2-IH^2} \ge \sqrt{R^2-IA^2} \Rightarrow r \geq \sqrt{R^2-IA^2}$.\\
		Diện tích của hình tròn $(C)$ nhỏ nhất khi bán kính $r$ nhỏ nhất, tức là $r=\sqrt{R^2-IA^2} \Leftrightarrow H\equiv A$.\\
		Khi đó $IA\perp(P) \Rightarrow$ mặt phẳng $(P)$ nhận $IA=(-1;-2;-1)$ làm một véctơ pháp tuyến.\\
		Vậy phương trình mặt phẳng $(P)\colon -x-2y-(z-2)=0\Leftrightarrow x+2y+z-2=0$.
	}
\end{ex}
%%%%==============EX_7============%%%
\begin{ex}%[2H5C3-3]
	Trong không gian $Oxyz$ cho mặt cầu $(S)\colon (x-1)^2+(y+2)^2+(z-3)^2=27$. Gọi $(\alpha)$ là mặt phẳng đi qua $2$ điểm $A(0;0;-4)$, $B(2;0;0)$ và cắt $(S)$ theo giao tuyến là đường tròn $(C)$ sao cho khối nón có đỉnh là tâm của $(S)$, là hình tròn $(C)$ có thể tích lớn nhất. Biết mặt phẳng $(\alpha)$ có phương trình dạng $a x+b y-z+c=0$, khi đó $a-b+c$ bằng
	\choice
	{$8$}
	{$0$}
	{$2$}
	{\True $-4$}
		\loigiai{
		$\bullet$ Vì $(\alpha)$ qua $A$ ta có $-(-4)+c=0\Rightarrow c=-4$.\\
		$\bullet$ Vì $(\alpha)$ qua $B$ ta có $2a+c=0\Rightarrow a=2$.\\
		$\Rightarrow(\alpha)\colon 2x+b y-z-4=0$.\\
		$\bullet$ Mặt cầu $(S)$ có tâm $I(1;-2; 3), R=3\sqrt{3}$.\\
		$\bullet$ Chiều cao khối nón $h=d_{(I, \alpha)}=\dfrac{|2-2b-3-4|}{\sqrt{4+b^2+1}}=\dfrac{|2b+5|}{\sqrt{b^2+5}}$.\\
		$\bullet$ Bán kính đường tròn $r=\sqrt{R^2-h^2}=\sqrt{27-\left(\dfrac{|2b+5|}{\sqrt{b^2+5}}\right)^2}=\sqrt{27-\dfrac{(2b+5)^2}{b^2+5}}$.\\
		$\bullet$ Thể tích khối nón $V=\dfrac{1}{3} \pi r^2 h=\dfrac{1}{3} \pi\left(27-\dfrac{(2b+5)^2}{b^2+5}\right) \dfrac{|2b+5|}{\sqrt{b^2+5}}$.\\
		$\bullet$ Tới đây ta có thể thử các trường hợp đáp án. Hoặc ta làm tự luận như sau\\
		Đặt $t=\dfrac{|2b+5|}{\sqrt{b^2+5}}$ và xét hàm số $f(t)=\left(27-t^2\right) t$ trên đoạn $[0; 3\sqrt{3}]$.\\
		Ta có $f'(t)=27-3t^2$.\\ $f'(t)=0\Leftrightarrow\hoac{&t=3\\ &t=-3~(\text{loại}).}$ \\
		Ta có bảng biến thiên
		\begin{center}
			\begin{tikzpicture}
				\tkzTabInit[nocadre=false,lgt=2,espcl=3,deltacl=0.5]
				{$x$/1 , $f'(x)$/1 , $f(x)$/2}
				{$0$,$3$,$3\sqrt{3}$}
				\tkzTabLine{,+,0,-,}
				\tkzTabVar{-/$0$,+/$54$,-/$0$,}
			\end{tikzpicture}
		\end{center}
		Do đó thể tích khối nón lớn nhất khi và chỉ khi\\ $t=3\Leftrightarrow\left(\dfrac{|2b+5|}{\sqrt{b^2+5}}\right)^2=3^2 \Leftrightarrow 4b^2+20b+25=9b^2+45$
		$\Leftrightarrow 5b^2-20b+20=0\Leftrightarrow b=2$.\\
		Vậy $a-b+c=-4$.
	}
\end{ex}
%%%%==============EX_8============%%%
\begin{ex}%[2H5C3-3]
	Trong không gian với hệ trục tọa độ $Oxyz$, cho mặt cầu $(S)$ có phương trình là $x^2+y^2+z^2-2x-2y-6z+7=0$. Cho ba điểm $A$, $M$, $B$ nằm trên mặt cầu $(S)$ sao cho $AMB=90^{\circ}$. Diện tích tam giác $AMB$ có giá trị lớn nhất bằng?
	\choice
	{\True $4$}
	{$2$}
	{$4\pi$}
	{Không tồn tại}
		\loigiai{
		Ta có $(S)\colon (x-1)^2+(y-1)^2+(z-3)^2=4\Rightarrow(S)$ có tâm $I(1;1;3)$ và bán kính $R=2$.\\
		Từ $A$, $M$, $B$ nằm trên mặt cầu $(S)$ và $\widehat{AMB}=90^{\circ} \Rightarrow AB$ qua $I \Rightarrow AB=2R=4$.\\
		Ta có $S_{AMB}=\dfrac{1}{2}\cdot MA \cdot MB \leq \dfrac{MA^2+MB^2}{4}=\dfrac{AB^2}{4}=4$.\\
		Dấu \lq\lq=\rq\rq xảy ra khi $ MA=MB=\dfrac{AB}{\sqrt{2}}=2\sqrt{2}$ và $AB=4$.\\
		Do đó diện tích tam giác $AMB$ có giá trị lớn nhất bằng $4$.
	}
\end{ex}
%%%%==============EX_9============%%%
\begin{ex}%[2H5C3-3]
	Trong không gian $Oxyz$, cho điểm $M\left(\dfrac{1}{2};\dfrac{\sqrt{3}}{2}; 0\right)$ và mặt cầu $(S)\colon x^2+y^2+z^2=8$. Đường thẳng $d$ thay đổi, đi qua điểm $M$, cắt mặt cầu $(S)$ tại hai điểm phân biệt $A$, $B$. Tính diện tích lớn nhất $S$ của tam giác $OAB$.
	\choice
	{\True $S=\sqrt{7}$}
	{$S=4$}
	{$S=2\sqrt{7}$}
	{$S=2\sqrt{2}$}
		\loigiai{
		Mặt cầu $(S)$ có tâm $O=(0;0;0)$, bán kính $R=2\sqrt{2}$.\\
		Ta có $OM=1\Rightarrow M$ nằm trong mặt cầu. \\
		Gọi $I$ là trung điểm $AB \Rightarrow OI\perp AB$.\\
		Đặt $x=OI\leq OM\Rightarrow 0<x\le 1$.\\
		Khi đó $S_{\triangle OAB}=\dfrac{1}{2}\cdot OI\cdot AB=OI \sqrt{R^2-O I^2}=x \sqrt{8-x^2}=f(x)$.\\
		$\Rightarrow f'(x)=\dfrac{2\left(4-x^2\right)}{\sqrt{8-x^2}}(0< x \le 1)$. \\
	Bảng biến thiên
\begin{center}
	\begin{tikzpicture}
		\tkzTabInit[nocadre=false,lgt=2,espcl=5,deltacl=0.5]
		{$x$/1 , $f'(x)$/1 , $f(x)$/2}
		{$0$,$1$}
		\tkzTabLine{,+,}
		\tkzTabVar{-/$0$,+/$\sqrt{7}$}
	\end{tikzpicture}
\end{center}
		Vậy $\max S_{\triangle O AB}=\sqrt{7}$ khi $OI=1$ hay $I\equiv M$.
	}
\end{ex}
%%%%==============EX_10============%%%
\begin{ex}%[2H5C3-3]
	Trong không gian với hệ trục $Oxyz$ cho hai đường thẳng $\Delta_1\colon \dfrac{x+1}{2}=\dfrac{y+1}{1}=\dfrac{z+1}{2}$ và $\Delta_2\colon \dfrac{x-1}{2}=\dfrac{y-1}{2}=\dfrac{z-1}{1}$. Tính diện tích mặt cầu có bán kính nhỏ nhất, đồng thời tiếp xúc với cả hai đường thẳng $\Delta_1$ và $\Delta_2$.
	\choice
	{$R=\dfrac{\sqrt{17}}{2}$}
	{$R=\dfrac{\sqrt{17}}{3}$}
	{$R=\dfrac{\sqrt{17}}{6}$}
	{\True $R=\dfrac{\sqrt{17}}{17}$}
		\loigiai{
		Gọi $A$, $B$ là hai điểm thuộc lần lượt $\Delta_1$ và $\Delta_2$ sao cho $AB$ là đoạn thẳng vuông góc chung giữa $2$ đường.\\
		Gọi $M$ là trung điểm $AB$. Dễ có mặt cầu tâm $M$ bán kính $R=\dfrac{AB}{2}$ tiếp xúc với hai đường thẳng $\Delta_1$ và $\Delta_2$ là mặt cầu có bán kính bé nhất.\\
		Ta có tọa độ theo tham số của $A$, $B$ lần lượt là
		$A\left(2t_1-1;t_1-1;2t_1-1\right)$ và $B\left(2 t_2+1; 2 t_2+1; t_2+1\right)\\ \Rightarrow \overrightarrow{AB}\left(2 t_2-2 t_1+2; 2 t_2-t_1+2; t_2-2 t_1+2\right)$.\\
		Có $\overrightarrow{u_1}(2;1;2)$ và $\overrightarrow{u_2}(2;2;1)$ lần lượt là $2$ vectơ chỉ phương của $\Delta_1$ và $\Delta_2$ nên
		$\heva{&\overrightarrow{AB} \perp \overrightarrow{u_1} \\
			&\overrightarrow{AB} \perp \overrightarrow{u_2}.}$\\
		$\Leftrightarrow\heva{&\left(2t_2-2t_1+2\right) \cdot 2+\left(2t_2-t_1+2\right) \cdot 1+\left(t_2-2t_1+2\right) \cdot 2=0\\
			&\left(2t_2-2t_1+2\right) \cdot 2+\left(2t_2-t_1+2\right) \cdot 2+\left(t_2-2t_1+2\right) \cdot 1=0}$\\
		$\Leftrightarrow\heva{&8t_2-9t_1+10=0\\
			&9t_2-8t_1+10=0} \Leftrightarrow\heva{&t_1=\dfrac{10}{17} \\
			&t_2=\dfrac{-10}{17}.}$\\
		$ \Rightarrow A\left(\dfrac{3}{17}; \dfrac{-7}{17}; \dfrac{3}{17}\right); B\left(\dfrac{-3}{17}; \dfrac{-3}{17}; \dfrac{7}{17}\right); \overrightarrow{AB}\left(\dfrac{-6}{17}; \dfrac{4}{17}; \dfrac{4}{17}\right)$.\\
		$R=\dfrac{AB}{2}=\dfrac{1}{2} \cdot \dfrac{\sqrt{(-6)^2+4^2+4^2}}{17}=\dfrac{\sqrt{17}}{17}$.\\
		Bán kính mặt cầu cần tính là $R=\dfrac{\sqrt{17}}{17}$.
	}
\end{ex}